\chapter{The checks, and the checks on the checks}

Enforcement infrastructure has a failure mode of its own: a check that cannot
pass its own rule gets an exemption, and an exemption is where enforcement starts
to rot. Three entries below are about checks that had to be built so that they
could survive themselves.

\begin{defect}{The dash checker was its own first violation}
\symptom{The first draft of the checker that forbids two specific characters
declared those two characters as string literals, which made the file an
immediate violation of the rule it enforces.}
\diagnosis{A checker that names what it forbids will contain what it forbids.}
\resolution{Build the characters from their codepoints at import.}
\instead{Exempting the checker from its own scan. The same reasoning applies to
the commit message checker, whose blocklist would otherwise have been the one
file in the repository carrying every string the repository forbids; its patterns
are stored encoded and decoded at import. Two tests do the same thing for the
same reason. This is a small point and it is a real one.}
\end{defect}

\begin{defect}{A traceability check that states its rule positively}
\symptom{The obvious design flags every numeric literal in prose and then excuses
most of them.}
\diagnosis{The exclusion list grows until the check means nothing. That is not a
hypothetical: it is the failure mode of essentially every lint rule that starts
by flagging everything.}
\resolution{State the rule positively. Three shapes count as a measurement: a
percentage, a probability shaped decimal, and a number sharing a line with a
measurement word. Years, dates, version strings, reference numbers, list markers,
code blocks and addresses are never flagged at all. A line is cleared either by a
result file reference or by an explicit annotation carrying a written reason, so
that waving the check through leaves a written trace with a name on it.}
\instead{Flag everything and maintain an exclusion list. Chosen against
deliberately, and the reasoning is in the checker's own docstring so that the
next person to be annoyed by a false negative knows what the trade was.}
\end{defect}

\begin{defect}{The traceability check passed on citations of files that did not
exist}
\symptom{Until the first result files existed the check could only look for the
\emph{shape} of a citation, and a check that matches a shape passes on a citation
of a file nobody committed.}
\diagnosis{That is the same thing as no citation at all, with more characters.}
\resolution{A resolve mode that walks the whole chain for every reference: the
path must exist, a manifest must sit beside it or above it, the manifest must
name a commit this repository can produce an object for, and the manifest must
not mark the run dirty. The scan and the resolve are separate so that the
scanner's own tests can run over strings with no repository behind them.}
\instead{Leaving it as a shape match. It had been green since the first phase and
would have stayed green forever, because it never asked the version control
system for anything.}
\end{defect}

\begin{defect}{The first genuinely red gate, and it was correct}
\symptom{The traceability job went red on a push, on a check that had passed
locally minutes earlier, reporting that a manifest's commit does not resolve in
this repository.}
\diagnosis{The job checked out at the default depth of one, so the clone held
only the tip commit and every citation of a run taken in an earlier commit
failed. The message was true of the clone and false of the repository.}
\resolution{Full history on that job, the same setting the ancestry job already
carried for the same reason: both checks ask questions about history, and history
is the input.}
\instead{Fixing it from the error message without reproducing it. The failure was
reproduced locally with a shallow clone before anything was changed, because a
failure that is fixed without being reproduced is a failure that is guessed at.

Worth recording for a second reason. It is the first time a gate in this project
went red on a real defect rather than on a deliberately broken branch, so the
argument for proving that each gate can fail now has a natural example beside its
manufactured ones.}
\end{defect}

\begin{defect}{The proof of failure exercise found two real problems}
\symptom{Each continuous integration job was deliberately broken on a scratch
branch, one at a time, to prove it can go red. Two of the runs found something
that was not the intended breakage.}
\diagnosis{The packaging job's first attempt assumed the installer would place
the console script at a fixed path, which is a guess and was wrong on the runner.
Separately, a deliberate taxonomy breakage failed both the reachability job and
the test job.}
\resolution{The packaging job now sets its own installer home and binary
directory, which both fixes the assumption and gives the built artefact a
genuinely clean environment to install into. The double failure was left alone,
because it is the correct outcome: the property is asserted in two independent
places on purpose, and a defect that only one of them noticed would mean the other
was not doing its job.}
\instead{Skipping the exercise on the grounds that all six jobs were green. A
gate that has only ever been green is indistinguishable from a gate that always
returns green, and the run links for each observed failure are recorded in the
repository so that the claim is checkable rather than asserted.}
\end{defect}

\begin{defect}{Placeholder modules that sank the coverage gate}
\symptom{Measured coverage of the library sat well below the gate at a point when
the library was a taxonomy, a small command line surface, and a set of empty
module placeholders.}
\diagnosis{The placeholders were written with a single future import each. That
is one executable statement per file, thirty of them, none ever imported.}
\resolution{Strip the import. The placeholders then have zero statements, which
is what the coverage gate should see: the phases that fill them bring their own
tests.}
\instead{Lowering the gate for the first phase, or excluding the placeholders in
the coverage configuration. Both would have needed to be undone later and one of
them would have been forgotten.}
\end{defect}

\begin{defect}{A reachability check that does more than it was asked to}
\symptom{Not a defect. A deliberate strengthening, recorded because it has a
sharp edge.}
\diagnosis{The rule that the taxonomy has exactly one definition in code is
enforceable mechanically by scanning the source and the scripts for label
literals written anywhere else.}
\resolution{It does that. The scan is limited to tokens that cannot be confused
with ordinary prose, meaning the hyphenated specification tokens and the
uppercase extras.}
\instead{Scanning for every token. Several specification tokens are also ordinary
English words, and a check that fires on one of them in a comment is a check
somebody turns off. A check that cries wolf is worse than no check, because it
trains its reader to ignore it.}
\end{defect}

\begin{defect}{A first remote that was abandoned}
\symptom{The repository was initially pushed to a remote whose name differed in
case and separator from the one intended.}
\diagnosis{A naming mistake, caught before any run link had been recorded
anywhere.}
\resolution{The remote was repointed before anything referenced it, so no
recorded continuous integration link refers to the abandoned repository.}
\instead{Deleting the abandoned remote. The credentials in use do not carry
delete scope, so it is left in place and recorded here rather than quietly
forgotten.}
\end{defect}
